\chapter{Modelo de automatización e infraestructura como código}
\label{chap:modelo_automatizacion}

\section{Modelo de automatización e infraestructura como código}
\label{sec:modelo_automatizacion_iac}

\subsection{Fundamentos del modelo adoptado}

El repositorio adopta un modelo de infraestructura como código orientado a control de estado, no a ejecución aislada de comandos. La plataforma se describe mediante artefactos versionados que combinan inventario, roles, variables, plantillas y políticas de validación, de modo que la intención técnica y su implementación quedan en el mismo espacio de gobernanza. Esta decisión permite que cada cambio sea auditable en términos de alcance, impacto y coherencia con el resto del sistema, en lugar de depender de intervenciones manuales no trazables entre nodos.

En este marco, la automatización no se interpreta como una secuencia ad hoc de operaciones, sino como un proceso de convergencia sobre estado deseado. El proyecto evita rutas paralelas de provisión y concentra su lógica de despliegue en un flujo principal definido en código, lo que reduce ambigüedad operativa y desacopla la estabilidad del clúster del conocimiento tácito de una persona operadora. La reproducibilidad se deriva de esa propiedad: la misma definición declarativa conduce al mismo patrón de resultado cuando se aplica sobre el inventario vigente y sus variables efectivas.

La consistencia metodológica del capítulo se apoya en hechos concretos del repositorio. El diseño separa provisión, validación y limpieza destructiva en circuitos distintos: el aprovisionamiento ordinario se resuelve en \texttt{site.yml}, las comprobaciones se ejecutan en roles dedicados de validación y la remoción de componentes se encapsula en \texttt{cleanup\_slurm\_gpu.yml}. Esta frontera funcional evita que operaciones de naturaleza diferente compartan el mismo trayecto de ejecución y mantiene una semántica clara entre construir estado objetivo, verificar estado alcanzado y desmontar una instalación para reprovisión.

\subsection{Arquitectura de orquestación}

La orquestación se centraliza en un único entrypoint, \texttt{site.yml}, como decisión de diseño deliberada. Ese archivo estructura el despliegue por dominios técnicos en una secuencia de etapas encadenadas, donde cada bloque prepara precondiciones que el siguiente requiere para converger sin contradicciones. El valor de esta arquitectura no depende de fijar un número rígido de pasos, sino de preservar un orden lógico que estabiliza el sistema por capas y evita activar componentes sobre un contexto incompleto.

La progresión observable en el repositorio puede sintetizarse en macrofases funcionales. Esta agregación no reemplaza la granularidad de los plays, pero facilita entender cómo se articula la convergencia completa desde baseline hasta validación final sin perder correspondencia con roles reales.

\begin{table}[H]
\centering
\caption{Macrofases de orquestación y roles principales del repositorio}
\label{tab:modelo_orquestacion_macrofases}
\begin{tabular}{|p{0.34\textwidth}|p{0.58\textwidth}|}
\hline
\textbf{Macrofase} & \textbf{Roles principales involucrados} \\
\hline
Baseline del sistema y acceso & \texttt{common}, \texttt{users\_ssh} \\
\hline
Red y perímetro & \texttt{network\_internal}, \texttt{cluster\_routing}, \texttt{firewall} \\
\hline
Aceleración GPU & \texttt{nvidia\_cuda} \\
\hline
Almacenamiento compartido & \texttt{nfs\_hpc} \\
\hline
Base de datos de accounting & \texttt{mariadb\_server}, \texttt{slurm\_db\_prep} \\
\hline
Servicios de identidad y autenticación & \texttt{slurm\_identities}, \texttt{munge} \\
\hline
Scheduler (control y cómputo) & \texttt{slurm\_facts}, \texttt{slurm\_rpm\_build}, \texttt{slurm\_install}, \texttt{slurm\_controller}, \texttt{slurm\_compute} \\
\hline
Entorno de aplicación & \texttt{llm\_env} \\
\hline
Validación operativa & \texttt{validate}, \texttt{slurm\_validate} \\
\hline
\end{tabular}
\end{table}

La coherencia de esta arquitectura se refuerza porque el orden de macrofases respeta dependencias técnicas reales. La inicialización de sistema y acceso precede al endurecimiento de red y firewall; posteriormente se habilitan servicios de datos, identidad y scheduler; y la validación queda desacoplada como una capa de verificación posterior al cambio. Esta secuencia disminuye la probabilidad de estados intermedios inconsistentes y acota los puntos de fallo asociados a despliegues fuera de orden.

\subsection{Control del cambio y segmentación operativa}

El control del cambio se implementa combinando secuenciación por etapas, segmentación por grupos de hosts y uso estratégico de tags. \texttt{site.yml} aplica dominios de configuración sobre alcances explícitos como \texttt{all}, \texttt{hpc\_master}, \texttt{all:!hpc\_master}, \texttt{slurm\_all} y \texttt{slurm\_compute}, lo que evita ambigüedad sobre dónde debe converger cada fase. Esta segmentación reduce acoplamiento entre planos funcionales y permite que un mismo modelo de orquestación preserve diferencias de rol entre nodos sin duplicar lógica declarativa.

Las tags no actúan como metadato ornamental, sino como mecanismo de recorte controlado del cambio. Al estar alineadas con dominios técnicos del repositorio, permiten delimitar intervenciones sin abandonar la trayectoria central de \texttt{site.yml}. Con ello, la granularidad de ejecución se mantiene dentro de un marco coherente, de manera que los ajustes parciales continúan sujetos a la misma jerarquía de variables, al mismo inventario y al mismo control histórico del código.

La reducción de riesgo se completa al mantener el playbook de limpieza fuera del flujo principal. \texttt{cleanup\_slurm\_gpu.yml} encapsula operaciones destructivas en una ruta separada, con enfoque explícito de seguridad operacional, incluyendo ejecución serial y preservación de conectividad administrativa. Esta separación impide que tareas de desmontaje compitan semánticamente con la provisión ordinaria y protege la operación regular frente a activaciones accidentales de remoción.

\subsection{Gestión jerárquica de variables y secretos}

La configuración del clúster se gobierna mediante una jerarquía explícita de variables que refleja alcance y especialización. El plano global se concentra en \texttt{group\_vars/all/vars.yml}; la especialización del controlador se define en \texttt{group\_vars/hpc\_master.yml}; los ajustes puntuales se ubican en \texttt{host\_vars}; y cada rol mantiene una base sobreescribible en \texttt{defaults}. Esta estructura permite adaptar comportamiento por dominio sin romper coherencia global y minimiza hardcodes dentro de tareas.

La precedencia entre capas transforma la configuración en un sistema determinista de resolución de estado. Primero se establece la intención transversal, luego se especializa por grupo y finalmente se ajusta por host cuando la topología o los recursos lo requieren. En términos de control metodológico, esta organización facilita auditar por qué un nodo converge de cierta forma y en qué nivel de la jerarquía se decidió ese comportamiento.

La gestión de secretos se integra en el mismo marco sin mezclar valores sensibles con lógica operativa visible. El repositorio utiliza \texttt{group\_vars/all/vault.yml} cifrado, y las variables de configuración referencian claves \texttt{vault\_*}, como ocurre con parámetros de base de datos y credenciales de elevación en segmentos específicos del inventario. Esta práctica preserva confidencialidad en almacenamiento, mantiene ejecutabilidad del flujo declarativo y conserva trazabilidad de cambios en torno a secretos sin exponerlos en texto plano.

\subsection{Consistencia distribuida e idempotencia}

La consistencia distribuida del modelo se sostiene en la convergencia repetible de roles, no en sincronización manual posterior. La recolección sistemática de facts en etapas de \texttt{site.yml}, junto con la separación por grupos y variables de alcance, permite que cada nodo reciba configuración acorde a su función dentro de un patrón común de estado. En servicios con interacción entre nodos, esta regularidad reduce deriva de configuración y favorece comportamiento homogéneo del conjunto.

La idempotencia se implementa mediante uso de módulos declarativos y control explícito del cambio efectivo. La lógica del repositorio prioriza estados observables de paquetes, servicios y archivos, y reserva los reinicios para handlers cuando corresponde, evitando reinicios indiscriminados en cada ejecución. En la capa de verificación, los roles \texttt{validate} y \texttt{slurm\_validate} se mantienen desacoplados de la provisión y emplean tareas de lectura con criterios de no cambio, de modo que la auditoría operativa no altera el estado que pretende medir.

Esta combinación de convergencia declarativa, control de efectos y validación separada convierte la repetición de ejecución en una prueba de estabilidad del modelo. Si el estado objetivo ya fue alcanzado, la automatización tiende a mantenerlo en lugar de reconfigurarlo innecesariamente; si existen desvíos, los corrige dentro del mismo marco de responsabilidades. El resultado es un ciclo de operación técnicamente coherente en el que cambio, verificación y mantenimiento comparten una base metodológica única y verificable.
